\documentclass[11pt]{letter}

\usepackage[T1]{fontenc}
\usepackage[utf8]{inputenc}
\usepackage[a4paper,margin=25mm]{geometry}
\usepackage{microtype}
\usepackage[colorlinks=true,urlcolor=blue,linkcolor=black,citecolor=black]{hyperref}

\address{%
Javier Macias-Guarasa\\
Departamento de Electronica\\
Universidad de Alcala\\
Alcala de Henares, Spain\\
\href{mailto:javier.maciasguarasa@uah.es}{javier.maciasguarasa@uah.es}}

\signature{Javier Macias-Guarasa\\
Corresponding author, on behalf of all authors}

\date{\today}

\begin{document}

\begin{letter}{Editors\\\textit{Scientific Data}}

\opening{Dear Editors of \textit{Scientific Data},}

We are pleased to submit the manuscript entitled
\textbf{``A Distributed Acoustic Sensing Dataset for Vessel Detection and
Localization in Submarine Cable Protection''}, in which we describe the
Marlinks North Sea Distributed Acoustic Sensing Dataset (Marlinks-NS DAS), for
consideration as a Data Descriptor in \textit{Scientific Data}.

The context motivating this dataset is the need to monitor submarine
telecommunications and power cables, which are critical infrastructure exposed
to accidental damage and deliberate interference. Distributed acoustic sensing
(DAS) can support this objective by turning optical-fibre cables into dense
arrays of vibration sensors. However, openly reusable and labelled DAS datasets
acquired in submarine environments, particularly those synchronised with vessel
information, remain scarce. The Marlinks-NS DAS dataset addresses this data gap
through a ten-day acquisition campaign in the North Sea. It contains 74,771
labelled 10-s observations, each comprising spectral-energy features from 250
sensing channels and 100 frequency bands, together with AIS-derived
vessel-to-cable distance labels and anonymised vessel type, length, and beam
information. The dataset supports research on vessel detection and
vessel-to-cable distance estimation for submarine-infrastructure monitoring and
protection.

We consider the manuscript well suited to \textit{Scientific Data} because it
is centred on the creation, curation, organisation, validation, and reuse of a
publicly deposited research dataset. It provides a detailed account of the data
curation process, including DAS acquisition and preprocessing, AIS trajectory
interpolation and synchronisation, geospatial distance-label generation,
spectral-feature extraction, and the organisation of the released HDF5 data.
The deposited research object also includes documentation, supporting metadata,
frequency-band definitions, daily environmental summaries, and archived Python
examples for inspecting, loading, checking, and partitioning the data. The
technical validation goes beyond baseline classification and regression
experiments by also examining physical-passage overlap and vessel-identity
recurrence across the recommended day-wise folds, as well as environmental
variability across those folds. These analyses are intended to help prospective
users assess dataset quality, acquisition diversity, and the suitability of the
recommended evaluation protocol.

The dataset may be useful to researchers developing machine-learning methods
for distributed acoustic sensing, studying vessel detection and
vessel-to-cable distance estimation for submarine-infrastructure monitoring and
protection, investigating marine acoustics and submarine-fibre sensing, and
developing spatial, spectral, or spatiotemporal learning methods for distributed
sensor arrays. The dataset, documentation, metadata, and example code are openly
available on Zenodo at
\href{https://doi.org/10.5281/zenodo.15611778}{https://doi.org/10.5281/zenodo.15611778}.

For transparency, the DAS and AIS measurements underlying this dataset were
also used in a related application-centred article that was recently accepted
for publication in the \textit{IEEE Journal of Selected Topics in Applied Earth
Observations and Remote Sensing}. An early-access version is available at
\href{https://doi.org/10.1109/JSTARS.2026.3716768}{https://doi.org/10.1109/JSTARS.2026.3716768},
with minor editorial corrections currently being processed. That article
focuses on algorithm development and application performance. The present
manuscript explicitly describes the relationship to that publication and adds
substantial new and more detailed dataset-centred information to establish the
released data as an independently reusable research resource. A preprint of the
current Data Descriptor is also available on arXiv at
\href{https://doi.org/10.48550/arXiv.2607.28306}{https://doi.org/10.48550/arXiv.2607.28306}.

The manuscript is not under consideration by another journal, and all authors
have reviewed and approved both the manuscript and its submission to
\textit{Scientific Data}.

Thank you for considering our submission.

\closing{Sincerely,}

\end{letter}
\end{document}
